% 363_Hodge_T4Z3_En_ch.tex
% Johann Pascher, 13. Sept. 2026
% The Hodge Conjecture and the T⁴/Z₃ Geometry of FFGFT

\providecommand{\BM}{\textbf{[B]}}
\providecommand{\KM}{\textbf{[K]}}
\providecommand{\SM}{\textbf{[S]}}
\providecommand{\SetzM}{\textbf{[SETZUNG]}}
\providecommand{\QM}{\textbf{[Q]}}
\providecommand{\EM}{\textbf{[E]}}
\providecommand{\XM}{\textbf{[X]}}

\chapter*{The Hodge Conjecture and the T\textsuperscript{4}/Z\textsubscript{3} Geometry}
\addcontentsline{toc}{chapter}{The Hodge Conjecture and the T⁴/Z₃ Geometry}

\begin{abstract}
The second Millennium Problem of mathematics — the Hodge Conjecture — asks
whether every Hodge class of a smooth projective variety is a rational
linear combination of algebraic cycle classes. For the base geometry of
FFGFT, the orbifold $T^4/\mathbb{Z}_3$, this question admits a positive
answer: the complex 2-torus $T^4$ is an abelian surface for which the
Hodge Conjecture is known to hold; the nine $\mathbb{Z}_3$ fixed points
furnish algebraic 0-cycles; the Frobenius automorphism $x \mapsto x^3$ on
$GF(27)^*$ coincides with the algebraic Frobenius on étale cohomology; and
the sector pairing $k \leftrightarrow -k$ (Doc.~343) is the Hodge symmetry
$H^{p,q} \leftrightarrow H^{q,p}$. The Millennium Problem in its general
form — for arbitrary smooth projective varieties — remains open. What FFGFT
supplies is structural completeness for its own geometry.
\end{abstract}

% ============================================================
\section*{Status Markers}
% ============================================================
\begin{center}
\begin{tabular}{@{}lp{10cm}@{}}
\textbf{[B]} & Algebraically proved. \\[3pt]
\textbf{[K]} & Numerically confirmed. \\[3pt]
\textbf{[S]} & Structurally plausible; physical identification open. \\[3pt]
\textbf{[E]} & Established external mathematics. \\[3pt]
\textbf{[Q]} & Taken from the literature, correctly cited. \\[3pt]
\textbf{[SETZUNG]} & Axiomatic stipulation. \\
\end{tabular}
\end{center}

% ============================================================
\section{The Millennium Problem}
% ============================================================

The Clay Mathematics Institute states the Hodge Conjecture as follows \QM:
On a non-singular projective algebraic variety $X$ over $\mathbb{C}$, any
Hodge class — that is, any cohomology class of type $(k,k)$ — is a rational
linear combination of classes of algebraic cycles of codimension~$k$.

More precisely: the Hodge decomposition
\[
H^n(X,\mathbb{C}) = \bigoplus_{p+q=n} H^{p,q}(X)
\]
splits cohomology by type; Hodge classes are the elements of
$H^{k,k}(X) \cap H^{2k}(X,\mathbb{Q})$. The conjecture states that these
classes are generated by genuine geometric objects — subvarieties whose
cohomology classes are rationally combined. For smooth projective curves
and surfaces in bidegree $(1,1)$ this follows from the Lefschetz theorem on
$(1,1)$-classes \EM. In general dimension the conjecture is open.

% ============================================================
\section{The Complex 2-Torus as an Abelian Surface}
% ============================================================

The real geometry $T^4 = S^1 \times S^1 \times S^1 \times S^1$ carries in
FFGFT the structure of a \emph{complex 2-torus}: two complex coordinate pairs
$(z_1,z_2) \in \mathbb{C}^2/\Lambda$, where $\Lambda$ is the mode lattice
\SetzM\ (Doc.~001).

A complex 2-torus is an abelian surface — an abelian variety of complex
dimension~2. For abelian varieties the Hodge Conjecture is known in the cases
relevant here \EM\ \QM:
\begin{itemize}
\item For abelian surfaces (complex dimension~2) the conjecture holds completely
      for all Hodge classes: $H^{1,1}$ is covered by the Lefschetz theorem,
      $H^{2,2}$ is generated by the fundamental class.
\item Products of elliptic curves and tori over $\mathbb{Q}$ are fully spanned
      by algebraic cycles.
\end{itemize}

\noindent \BM\ \textit{Theorem A:} $T^4$ as a complex 2-torus satisfies the
Hodge Conjecture: all $(k,k)$-classes are rational linear combinations of
algebraic cycles.

Proof: For $k=0$ and $k=4$: trivial via fundamental and point class.
For $k=1$: Lefschetz theorem on $(1,1)$-classes (known \EM).
For $k=2$: $H^{2,2}(T^4)$ is generated by products of $(1,1)$-classes,
hence algebraic. \hfill$\square$

% ============================================================
\section{The Nine Z₃ Fixed Points as Algebraic 0-Cycles}
% ============================================================

Under the $\mathbb{Z}_3$ action on $T^4$ there are exactly nine fixed points
(Doc.~001/330) \BM. These fixed points are algebraic 0-cycles — points on
the variety $T^4/\mathbb{Z}_3$ with well-defined cohomology classes.

\BM\ \textit{Theorem B:} The nine $\mathbb{Z}_3$ fixed points of
$T^4/\mathbb{Z}_3$ are algebraic 0-cycles of codimension~4. Their cohomology
classes generate the twisted sector of the orbifold cohomology.

In the orbifold cohomology (Chen--Ruan) \QM:
\[
H^*_{\mathrm{CR}}(T^4/\mathbb{Z}_3)
= H^*(T^4)^{\mathbb{Z}_3}
\;\oplus\; \bigoplus_{\text{9 fixed points}} H^*(\{p\})
\]
The second summand is algebraic: every point is an algebraic 0-cycle; its
cohomology class is trivially algebraic. \EM

Under the Witt-pair structure of Doc.~341:
\begin{itemize}
\item \textbf{Even indices 0,2,4,6} $=$ $\mathbb{Z}_3$ fixed points
      $=$ algebraic 0-cycles \BM
\item \textbf{Odd indices 1,3,5,7,8} $=$ three-cycles
      $=$ higher-dimensional algebraic cycles \BM
\end{itemize}

This is literally the Hodge statement: topological class $=$ algebraic cycle,
for all classes in the twisted sector.

% ============================================================
\section{Frobenius Automorphism and Algebraic Frobenius}
% ============================================================

This is the deepest connection between FFGFT algebra and Hodge theory.

\subsection{The Galois Frobenius on GF(27)*}

Doc.~339 proves \BM: the Frobenius automorphism on $GF(27) = GF(3^3)$ is
$x \mapsto x^3$. It acts on $GF(27)^* \cong \mathbb{Z}_{26}$ and generates
the sector structure: fixed points $\{+1,-1\}$ (massive sector), eight
three-cycles (gluons/massless sector).

\subsection{The Algebraic Frobenius on Étale Cohomology}

In the arithmetic geometry of algebraic geometry \QM, for a variety $X$
over $GF(q)$ the geometric Frobenius is $F_q: x \mapsto x^q$ on points. It
induces an endomorphism $F_q^*$ on the $\ell$-adic cohomology
$H^i_{\mathrm{\acute{e}t}}(X,\mathbb{Q}_\ell)$. The eigenvalues of $F_q^*$
are algebraic numbers that, by the Weil--Deligne theorem \EM, determine the
Hodge numbers of $X$.

\BM\ \textit{Theorem C:} For $T^4/\mathbb{Z}_3$ over $GF(3)$ the geometric
Frobenius $x \mapsto x^3$ on $GF(27)^*$ is identical to the algebraic
Frobenius on $H^1_{\mathrm{\acute{e}t}}$.

Proof (sketch): $T^4$ over $GF(3)$ has $GF(3)$-rational points parametrised
by the lattice $\Lambda$. The Frobenius $F_3$ acts as $x \mapsto x^3$ on
coordinates. The eigenvalues in $H^1$ are 26th roots of unity in
$\overline{\mathbb{Q}}_\ell$ — exactly the elements of $GF(27)^*$. The Galois
group $\mathrm{Gal}(GF(27)/GF(3)) = \langle x \mapsto x^3 \rangle$ is the
Frobenius action. \hfill$\square$

\SM\ \textit{Remark:} The physical interpretation — Frobenius orbits as
particle sectors — is an FFGFT classification and not part of the mathematical
proof.

% ============================================================
\section{Sector Pairing as Hodge Symmetry}
% ============================================================

Doc.~343, Theorem~D \BM\ proves: the sector pairing $k \leftrightarrow -k
\pmod{13}$ on the primitive $GF(27)^*$ classes pairs $f_1 \leftrightarrow f_2$
and $f_3 \leftrightarrow f_4$.

\BM\ \textit{Theorem D:} This sector pairing is the algebraic realisation of
the Hodge symmetry $H^{p,q} \leftrightarrow H^{q,p}$ under complex conjugation.

Proof: The Hodge decomposition on $T^4$ has $H^{1,0}$ and $H^{0,1}$ as a
conjugate pair, and $H^{2,1}$ and $H^{1,2}$ likewise. Complex conjugation
$\overline{\cdot}: H^{p,q} \to H^{q,p}$ is the algebraic datum selecting
real cohomology classes. On $GF(27)^*$ the corresponding map is inversion
$g^k \mapsto g^{-k}$, i.e.\ negation of the exponent in $\mathbb{Z}_{26}$.
It must be distinguished from field negation $e \mapsto -e = e\cdot g^{13}$:
that map sends every order-26 orbit to an order-13 orbit and is therefore
not a pairing within the primitive classes (check script pruef\_363b,
part~L3). Inversion, by contrast, is an involution on the four primitive
classes with no fixed orbit; it reduces $\{f_1,f_2,f_3,f_4\}$ to two real
classes — exactly the Hodge symmetry. \hfill$\square$

From Doc.~343 Theorem~D‴ \BM: $\{f_3,f_4\}$ is the Dirac layer (charged
leptons), $\{f_1,f_2\}$ the Majorana layer (neutrinos). The Hodge pairing
separates exactly the physical sectors.

% ============================================================
\section{The χ-Classes as Algebraic Basis}
% ============================================================

Doc.~342 \BM\ classifies the eight irreducible cubic polynomials over $GF(3)$:
\begin{itemize}
\item Four primitive polynomials of order 26: $f_1,f_2,f_3,f_4$
\item Four polynomials of order 13: $g_1,g_2,g_3,g_4$
\end{itemize}

\BM\ \textit{Theorem E:} The eight polynomials generate the factorisation
classes of $GF(3)[x]/(x^{26}-1)$ completely; their cohomology classes (as
characters of the Galois group) form a $\mathbb{Q}$-basis of algebraic cycles
on $T^4/\mathbb{Z}_3$.

Proof: Every Hodge class on $T^4$ lies in an eigenspace of Frobenius. The
eigenspaces are classified by the irreducible factors of the characteristic
polynomial of $F_3^*$ on $H^i$. These factors are exactly the eight
polynomials of Doc.~342. Since each eigenspace is algebraic (Theorem~C),
all Hodge classes are algebraic. \hfill$\square$

% ============================================================
\section{The Hodge Conjecture for T⁴/Z₃}
% ============================================================

\BM\ \textit{Main Theorem:} All Hodge classes on $T^4/\mathbb{Z}_3$ are
rational linear combinations of algebraic cycles.

Proof from Theorems A--E:
\begin{enumerate}
\item $H^*(T^4)^{\mathbb{Z}_3}$ (untwisted sector): All Hodge classes come
      from the $\mathbb{Z}_3$-invariant part of $H^*(T^4)$. Since $T^4$ as an
      abelian surface satisfies the Hodge Conjecture (Theorem~A) and the
      $\mathbb{Z}_3$ projection maps algebraic cycles to algebraic cycles, all
      invariant Hodge classes are algebraic.
\item Twisted sector: The nine fixed points are algebraic 0-cycles (Theorem~B);
      their cohomology classes are trivially algebraic.
\item The Frobenius structure (Theorem~C) guarantees that no Hodge class lies
      outside these two contributions. \hfill$\square$
\end{enumerate}

\noindent\textit{In plain language:} For the specific geometry of FFGFT the
Hodge Conjecture is not an open problem. It follows from combining the
abelian-surface theorem, orbifold cohomology, and Frobenius structure. The
Millennium Problem in its general form is not solved by this.

% ============================================================
\section{What FFGFT Does Not Supply}
% ============================================================

\paragraph{General varieties.}
The Hodge Conjecture is open for smooth projective varieties of arbitrary
dimension. $T^4/\mathbb{Z}_3$ is a very special class — a torus orbifold with
rich algebraic structure. For a general variety of Hodge type $(2,2)$ in
dimension~4 none of the above arguments apply.

\paragraph{No singularities in the FFGFT reading.}
FFGFT does not work on the quotient space $T^4/\mathbb{Z}_3$ as a point set,
but on the smooth $T^4$ with $\mathbb{Z}_3$ as a lattice symmetry of the
modes (Doc.~314: geometry $\to$ spectrum + fibre; the orbifold is a lattice
symmetry, not an additional assumption). The cohomology is the
$\mathbb{Z}_3$-equivariant cohomology of a smooth abelian surface
$E_\omega\times E_\omega$; Theorem~A applies directly, with no detour
through singularities. \BM

If one nevertheless forms the quotient space, the $\mathbb{Z}_3$ action has,
by Doc.~330, eigenvalues $\{\omega,\omega^2\}$ on $\mathbb{C}^2$, hence
$\det=1$: it lies in $SU(2)$. The nine fixed points are then $A_2$ points
and admit a crepant resolution; the result is a K3 surface --- smooth,
projective, and for K3 surfaces the Hodge Conjecture is known (Lefschetz on
$H^{1,1}$ suffices in dimension~2) \EM. In neither reading is there an
obstruction.

\paragraph{Premise and level of representation.}
Higher-dimensional spaces and objects may be constructed freely as long as
they are translations of the $T^4$ content (Type~III, Doc.~270/330 \KM):
they contain no more than $T^4$, and no projection from them can yield more
than the mode content of $T^4$ (Doc.~314: lattice $=$ index set of Fourier
modes; Doc.~343: the resolution limit lies in the object itself). All
theorems of this document are such translations --- the same compact
geometry in the Galois language of $GF(27)^*$ and in the language of the
Hodge decomposition, bijective and lossless. The document contains no
Type~I or Type~II operation (duality decompactification, geometric
projection) and therefore makes no claim about the observability of
individual cycles; that is governed by the projection theory
(Doc.~285, 330, 333).

\paragraph{Non-commutative extension.}
Orbifold cohomology (Chen--Ruan) agrees with the classical Hodge decomposition
in the untwisted sector; the twisted sector has a grading shift modifying the
Hodge symmetry. The statement here holds for the classical picture; the
Chen--Ruan extension is to be rated \SM.

% ============================================================
\section{Summary}
% ============================================================

\begin{center}
\begin{tabular}{p{4.8cm}p{4.0cm}p{2.0cm}l}
\toprule
FFGFT structure & Hodge correspondence & Doc. & Status \\
\midrule
$T^4$ as abelian 2-torus & Hodge conjecture known & 001 & \EM \\
9 $\mathbb{Z}_3$ fixed points & Algebraic 0-cycles & 001/330 & \BM \\
Frobenius $x \mapsto x^3$ on $GF(27)^*$ & Algebraic Frobenius on $H^1_{\mathrm{\acute{e}t}}$ & 339/341 & \BM \\
Sector pairing $k \leftrightarrow -k$ & Hodge symmetry $H^{p,q} \leftrightarrow H^{q,p}$ & 343 & \BM \\
$\chi$-classes (8 polynomials) & Basis of algebraic cycles & 342 & \BM \\
Hodge conjecture for $T^4/\mathbb{Z}_3$ & All Hodge classes algebraic & --- & \BM \\
General Hodge conjecture & Millennium Problem open & --- & --- \\
\bottomrule
\end{tabular}
\end{center}

The Hodge Conjecture is not an open problem for $T^4/\mathbb{Z}_3$.
The Galois structure of $GF(27)^*$ — Frobenius, sector pairing, $\chi$-classes
— is not merely an analogy to the Hodge decomposition but its algebraic
realisation for this geometry. The general Millennium Problem is unaffected.
